\section{Introduction}
\label{sec:intro}

Reproduction numbers summarise transmission under specified conditions and are routinely estimated from surveillance counts \citep{fraser2007estimating, gostic2020practical}. Renewal methods relate incidence to past infectiousness through a generation-interval distribution, and the width of their credible intervals shrinks as the number of cases grows \citep{cori2013new, thompson2019improved}. Work on low incidence, reporting noise and timely detection has made the limits of period-specific inference explicit \citep{parag2021improved, parag2022quantifying, parag2025asymmetric}, and temporal aggregation and reporting delays separate the reproduction number of the transmission process from the apparent reproduction number of an observed series \citep{nash2023estimating, ogigittins2024simulation, eales2024differences}.

Precision guidance for renewal estimators has been stated in cases. \citet{cori2013new} tabulate the number of incident cases a time window needs for a chosen posterior coefficient of variation and recommend starting estimation once at least 12 cases have accrued. Under a gamma prior with shape $a_0$, that coefficient is $(a_0+I_W)^{-1/2}$, where $I_W$ is the number of cases in the window, so it is a function of the window count alone. It measures counting error for the average transmission level over the window. The estimate is also often read as the transmission level now or as a typical level for the population. For those targets, variation between periods is an additional error component that more cases do not remove. The criterion is therefore correct for one target and incomplete for the other two, and the distinction depends on which target the analyst intends.

The variance structure behind this distinction is classical. Variance-component models split the variance of a unit-level estimate into a within-unit sampling term and a between-unit term \citep{searle1992variance, cochran1977sampling}. Random-effects meta-analysis estimates the between-study component on top of within-study error \citep{dersimonian1986meta}, and small-area estimation and empirical-Bayes disease mapping weight each direct estimate against a common mean according to the ratio of the two terms \citep{fay1979estimates, clayton1987empirical, gelman1999maps}. Funnel plots for institutional comparison draw control limits from the counting term and add an overdispersion component when units differ beyond chance \citep{spiegelhalter2005funnel}. Population ecology makes the same split between demographic and environmental stochasticity \citep{engen1998demographic, lande2003stochastic}, as do stochastic epidemic models \citep{breto2009time, he2010plug}. In each setting the count at which the two terms are equal is where a shrinkage estimator places half its weight on the direct estimate.

Several results used below are standard and are cited where they are applied. The mean squared error decompositions follow from the law of total variance \citep{searle1992variance, casella2002statistical}. The efficiency of the count-to-scale ratio for a Poisson mean is the Cramér--Rao bound \citep{lehmann1998theory}, the delta-method expansion of a ratio estimator is textbook material \citep{casella2002statistical}, and the square-root allocation of surveillance effort is Neyman allocation \citep{cochran1977sampling}. Offspring overdispersion follows \citet{lloydsmith2005superspreading}, and its estimation from small samples \citet{lloydsmith2007maximum}. A variance that scales as a power of the mean is Taylor's law \citep{taylor1961aggregation}, the alternative to the affine form that Section~\ref{sec:diagnostics} tests. The distinction between the realised quantity of a finite population and the parameter of the process that generated it follows \citet{graubard2002inference}. While these mathematical decompositions apply classical principles of variance analysis and survey sampling, the operational estimand linking empirical surveillance dispersion to transmission versus reporting mechanisms, the formalization of the crossover count $m_\times$ as an actionable shrinkage balance point, and the known-truth simulation audit on the real design are the distinct contributions of this work.

The paper makes four contributions. Section~\ref{sec:model} sets out three inferential targets for a period-specific reproduction number and derives their exact mean squared errors in a branching model with negative binomial offspring and binomial reporting, together with the correction required when the denominator is the previous observed count. Section~\ref{sec:reporting} shows that, when the reporting probability fluctuates between periods, single-period second moments of the counts cannot separate reporting volatility from transmission variation, and that the realised-mean target then retains a reporting term. Section~\ref{sec:diagnostics} defines three cross-sectional dispersion diagnostics for surveillance counts, with reference values for the counting slope under independent Poisson and branching Poisson counting. The same section reports a known-truth simulation study on the real sampling design, including a grid over the share of variance carried by transmission, the coverage of dependence-aware intervals and simulated reference distributions for the season and held-out comparisons. Section~\ref{sec:application} applies the diagnostics to 1,748 state-weeks of US influenza hospital admissions.

The diagnostics address a practical gap. Routine hospital surveillance does not observe the variables of a branching experiment. Admissions follow infection after a delay, aggregate heterogeneous populations and are recorded on a calendar schedule, and their count scale is an outcome of the epidemic. Reporting probabilities may change for all infections in the same period \citep{azmon2014estimating, stoner2019hierarchical}, and separating process from observation variation in time series is a known difficulty \citep{knape2008estimability, augermethe2016state}. The cross-sectional spread of growth-rate or reproduction-number estimates has been read as variation in transmission or behaviour between places and periods without subtracting the part of that spread that depends on case counts \citep{burghardt2022unequal, biggerstaff2014estimates}, although per-unit standard errors make the subtraction possible \citep{ives2021estimating}. Regressing the squared relative deviation of an apparent growth statistic on inverse case count across units, and reading the intercept as transmission variance or an error floor and the slope as counting noise, applies the classical decomposition to such data. Whether those readings hold has to be checked against data with a known truth, and Section~\ref{sec:diagnostics} reports that check. Section~\ref{sec:discussion} discusses real-time use, reporting that depends on transmission and the choice of generation interval. Proofs, extensions and supplementary analyses, including a cross-regime check on COVID-19 counts from mainland China, are in the appendices.
